% =====================================================================
% LUCRAREA COMPLETA, un singur fisier, de dat copy-paste in Overleaf.
% Template-ul ACL cu textul nostru in el, de la abstract la Limitations.
% Generat pe 30 iulie 2026 din intro_preview.tex (partile verbatim ale
% template-ului si Sectiunile 1-2) plus paste_body.tex si
% paste_limitations.tex, ambele rescrise, vezi mai jos.
%
% Numerotarea: 1 Introduction, 2 Related Work, 3 Approach (3.5 What was
% tuned), 4 Results, 5 What Did Not Work, 6 Conclusions, apoi Limitations
% nenumerotata. Fara Acknowledgments si fara apendice, vezi abaterile 6 si 7.
%
% ABATERILE DE LA TEMPLATE, toate opt:
%
%   1. Corpul lucrarii. Abstractul template-ului e inlocuit cu al nostru,
%      \section{Introduction} cu Sectiunile 1-6, si \section*{Limitations}
%      cu a noastra. Sectiunile de manual ale template-ului (Engines,
%      Preamble, Document Body, BibTeX Files) sunt SCOASE, cu tabelele si
%      figurile lor de exemplu. Consecinta: nu mai e nevoie de pachetul
%      mwe, si dispar \nocite{Ando2005,...} si citarile la Gusfield:97.
%
%   2. Bibliografia. Blocul filecontents scrie custom.bib cu cele 7 intrari
%      ale template-ului PLUS cele 18 ale lucrarii. Cele 7 ramin necitate,
%      dar inofensive: fisierul merge lipit si intr-un proiect ACL care are
%      deja un custom.bib de exemplu.
%
%   3. Doua adaugiri de preambul, dupa \usepackage{graphicx}: booktabs,
%      pentru \toprule / \midrule / \bottomrule din cele 5 tabele, si
%      \graphicspath{{figures/}} pentru cele 2 figuri.
%
%   4. Blocul \pinlabel din intro_preview.tex e STERS. Fixa numerele
%      Sectiunilor 3-5, care acum exista; referintele inainte din
%      introducere se rezolva singure.
%
%   5. Plasarea floats-urilor, si patru \renewcommand de fractii de float in
%      preambul. In paste_body.tex toate cele 7 floats erau declarate in bloc
%      la sfarsitul subsectiunii 4.4, desi cinci dintre ele sunt referite
%      prima data inca in Sectiunea 3. LaTeX le acumula si dedica coloane
%      intregi doar lor, cu spatiu elastic intre figuri: in prima versiune,
%      coloana din dreapta a paginii 5 era doar cele doua figuri, cu ~1.2
%      inch de alb intre ele si niciun rand de text. Acum fiecare float e
%      declarat inaintea titlului care urmeaza primei lui referinte, iar
%      topfraction / textfraction / floatpagefraction / bottomfraction sunt
%      relaxate ca sa incapa text in aceeasi coloana cu doua floats.
%      Rezultat, verificat la compilare: nicio coloana numai-floats, fiecare
%      figura si fiecare tabel are text sub el. tab:fail e singurul lasat pe
%      loc, pentru ca discutia lui e in Sectiunea 5, unde e si declarat.
%      Cropul PNG-urilor a fost verificat si NU ajuta: matplotlib a strins
%      deja marginile, bounding box-ul real recupereaza 2%.
%
%   6. Curatare de spatiu, 30 iulie 2026. PDF-ul a ajuns de la 9 la 8
%      pagini, fara sa se piarda nicio cifra si nicio revendicare. Trei
%      schimbari:
%        - \section*{Acknowledgments} SCOS. formatting.md: "Do not include
%          this section in the review version." Obligatoriu, nu economie de
%          spatiu. Vezi comentariul de la locul lui, inaintea bibliografiei.
%        - Apendicele de exemplu al template-ului SCOS, fara inlocuitor.
%          Lucrarea nu are apendice.
%        - Cinci duplicari verbatim, scoase. Patru intre proza si captions:
%          fraza despre +0.84 din 3.5 (daggerul din captionul lui
%          tab:ablation o spune, lipit de cifra), aritmetica 40% din
%          captionul lui tab:ablation (4.1 o spune, e teza paragrafului),
%          clauza despre SHAP din 4.3 (captionul lui tab:signals o spune),
%          si trimiterea circulara din 4.2 catre 3.5. A cincea, in
%          Limitations: paragraful despre cele 14 exemple repeta verbatim
%          din 4.2 clauza "surprisal was never computed for them and only
%          two carry eye-tracking". Scoasa acolo, pastrata in 4.2.
%
%      DRUM PARCURS, de nu repetat: tab:pca a fost intre timp mutat in
%      apendice, ca sa elibereze ~1/4 de pagina de continut. A functionat
%      la limita de pagini, dar arata rau: singur dupa cele 18 referinte,
%      ocupa o pagina intreaga cu 14 rinduri de text. Adus inapoi in 3.3.
%      Ultimele trei rinduri de bibliografie, care mai treceau pe pagina 9,
%      au fost recuperate de a cincea duplicare, cea din Limitations.
%
%      Stare masurata: 8 pagini in PDF, continutul numarat (Introduction ->
%      Conclusions) se termina in prima cincime a paginii 7, adica ~6.1
%      pagini din cele 8 permise de MRL. Limitations si References nu se
%      numara.
%
%   7. Citarea RALS si titlul, 30 iulie 2026, dupa confirmarea
%      coordonatorului. Cifrele de mai jos sunt verificate in PDF-ul RALS.
%
%      TITLUL e pus. Autorii nu: \author{} e gol, iar anonimizarea o face
%      template-ul, nu noi. acl.sty pune \acl@anonymizetrue pentru optiunea
%      "review" si tipareste "Anonymous ACL submission" in locul lui
%      \@author, deci continutul lui \author nu se vede la review oricum.
%
%      DOUA INTRARI NOI in bibliografie, 20 in total: anghel-etal-2025-rals
%      si cristea-nisioi-2024-machine. Vezi comentariile de la ele pentru
%      cele trei corectii fata de Anthology.
%
%      CINCI AFIRMATII REPARATE, toate devenite false sau contradictorii in
%      momentul in care RALS e citata:
%        - 2.1, prima fraza: "Romanian is not among them" era fals, RALS
%          publica 3921 de mostre adnotate pentru romana;
%        - 2.1, final: "Romanian appears in none of these resources"
%          restrins explicit la resursele multilingve inventariate;
%        - 3.1: "the gold labels for the test set are withheld" era fals,
%          sunt publice sub CC BY-NC-SA 4.0. Acum scrie ca sunt publice si
%          ca noi nu le-am folosit, cu consecinta ca metrica nu se poate
%          descompune;
%        - Limitations: acelasi fals, reparat la fel, plus declaratia ca
%          limitarea e a noastra si se poate elimina;
%        - Introducere, P3: recunoaste acum ca datele exista si ca le lasam
%          deoparte deliberat. Altfel se citea ca "nu exista date".
%
%      TREI ADAUGIRI care intaresc lucrarea:
%        - 2.2: cristea-nisioi-2024-machine transfera complexitatea in
%          romana prin traducere in engleza; RALS raporteaza r ~ 0.3 si
%          CONCLUZIONEAZA ca transferul cross-lingual de complexitate nu e
%          viabil pentru romana. Rezultatul nostru e contraexemplul. Cel mai
%          bun cadru de diferentiere pe care il are lucrarea.
%        - 4.1: reperul supervizat, r = 0.71 +/- 0.05, Ridge, 10-fold CV
%          grupat, pe EXACT setul nostru. Marcat explicit ca necomparabil,
%          fara nicio revendicare de depasire.
%        - 3.3 (A): RALS masoara corelatia Zipf-complexitate la -0.53 pentru
%          romana fata de -0.74 pentru engleza, adnotate la fel. Sprijin
%          independent, al altcuiva, pentru corectia pe lema.
%
%      Efect masurat: 20 de referinte, continutul creste la ~6.5 pagini din
%      8, PDF-ul revine la 9 pagini pentru ca bibliografia creste. Pagina 9
%      are 62 de rinduri de referinte, deci nu e o pagina orfana.
%
%   8. Acest comentariu.
%
% CE ARE NEVOIE PROIECTUL OVERLEAF: acl.sty si acl_natbib.bst, care vin cu
% template-ul, plus un folder figures/ cu DOUA png-uri din report/figures/:
% beeswarm.png si fig_corr_matrix.png.
%
% ATENTIE, consecinta lui [overwrite]: custom.bib se rescrie la fiecare
% compilare, deci intrari noi se adauga in blocul de mai jos, nu editand
% custom.bib in Overleaf, unde s-ar pierde.
%
% Doua reguli verificate, de nu incalcat la adaugari in .bib: niciun
% caracter arond in comentarii, ca BibTeX il citeste ca inceput de intrare;
% si niciun caracter Unicode in intrari, de aceea "Vorbesti Romaneste" e
% scris cu {\c s} la masala-etal-2024-vorbesti.
%
% ---------------------------------------------------------------------
% CE S-A SCHIMBAT FATA DE paste_body.tex SI paste_limitations.tex
%
% Amandoua erau la revizia de pe 27 iulie ora 16:09, in timp ce
% introducerea ajunsese la revizia 11. Aduse acum la aceleasi conventii.
% Cifrele, citarile, tabelele si figurile nu s-au atins.
%
%   (a) Em-dash scos integral, 29 de locuri in body si 9 in Limitations.
%       Cele cinci "---" ramase sunt celule goale de tabel, in tab:ablation
%       si tab:fail, nu em-dash de proza.
%   (b) Expresia "the system" ca desemnare a lucrarii noastre, interzisa la
%       revizia 10 a introducerii, scoasa din toate cele 8 locuri.
%   (c) Cadrul de competitie scos din revendicari. Termenul consecvent e
%       acum "the public evaluation score", nu "leaderboard", iar sursa
%       datelor e "released with the data", nu "supplied by the shared
%       task". ATENTIE: divulgarea din 3.5 si din Limitations a rams
%       INTEGRALA. Cele cinci cantitati au fost alese pe scorul public si
%       asta trebuie sa se poata citi; s-a schimbat incadrarea, nu faptul.
%       Referinta din Related Work la "the ten languages of the MLSP shared
%       task" a rams: e shared task-ul altora, citat ca resursa.
%   (d) O afirmatie rams falsa dupa corectiile de la revizia 8, in 4.2:
%       "the only real Romanian supervision in existence for this task".
%       Datele romanesti de complexitate exista, noi ALEGEM sa nu le
%       folosim. Reformulata ca singurele etichete pe care le folosim.
%   (e) Titlul primului paragraf din Limitations spunea "Three weights",
%       dar corpul lui numara trei ponderi SI doua constante de calibrare.
%       Acum "Five quantities".
%
% FIGURA SCOASA: fig_confusion.png, matricea de confuzie pe engleza
% held-out. Motivul principal, verificat in make_report_figures.py:93-121:
% codul antreneaza pe CompLex EN train+trial si testeaza pe CompLex EN
% test, adica engleza -> engleza, monolingv, desi captionul o numea
% "In-source check", iar 3.2 defineste in-source ca transfer intre seturi
% de limbi sursa diferite. In plus matricea contine strict mai putina
% informatie decat r = 0.753 din titlul ei propriu, pentru ca se obtine
% digitizand ambele axe la tertilele din train. Numarul a fost recuperat
% in 3.2, langa 0.515, unde contrastul dintre cele doua cuantifica costul
% trecerii intre limbi.
%
% TABEL DE p-VALUE, cerut si respins, cu calculul facut: la n = 1587, |r|
% critic pentru p < 0.05 e 0.0492, iar toate cele zece celule ale matricii
% de corelatii il depasesc, cea mai mica fiind B x C = 0.07 cu p = 5.3e-3.
% Tabelul ar fi fost zece rinduri constante, ar fi sugerat ca perechea pe
% care lucrarea o numeste "near-independent" e semnificativ corelata (0.07
% trece la limita si de Bonferroni pe 10 teste, critic 0.0704), si ar fi
% avut precizie falsa, pentru ca cele 78 de cuvinte fara gaze primesc
% mediana inainte de corelatie (fill() din solutie_finala.py:140), ceea ce
% umfla df-ul efectiv pe toate rindurile cu D. In loc de tabel, captionul
% lui fig:corr declara acum n si muta cititorul pe varianta comuna: 83%
% intre A si frecventa, 0.5% intre B si C.
%
% ---------------------------------------------------------------------
% DE FACUT INAINTE DE SUBMISIE (MRL @ EMNLP 2026, 10 august, double-blind)
%
%   Titlul si provenienta datelor sunt REZOLVATE, vezi abaterea 8.
%
%   1. tab:fail, rindul Lasso, in-source 0.569. Singura valoare in-source
%      fara log salvat. analysis_lasso.py afiseaza ambele metrici, deci o
%      rulare confirma daca e transfer_avg (cum e trecut) sau EN->ML.
%   2. r SI R^2 PE SETUL COMPLET: REZOLVAT PE 30 IULIE 2026, SUNT IN TEXT.
%      Etichetele RoLCP sunt publice (github.com/senisioi/RALS, data.zip,
%      parola "ralsresources" e in README-ul lor, CC BY-NC-SA 4.0).
%      Scorat submission.csv pe ele, join 1587/1587 pe cheie completa
%      (text_name, token, text), deci setul nostru de test E RoLCP exact:
%        r      = 0.7770   CI 95% [0.7567, 0.7958]
%        R^2    = 0.6028
%        compozit recalculat = 68.99  (returnat 69.74, diferenta 0.75)
%        gold: medie 0.2701, std 0.1674
%      Scrise in TREI locuri: 2.1 (paragraf nou "Our result on the same
%      set", langa cifra lor), 3.1 (finalul paragrafului de metrica) si 4.1
%      (paragraful reperului, rescris). In toate trei se spune ca etichetele
%      au fost citite doar dupa inghetarea predictorului, deci e raportare,
%      nu tuning. Doua paragrafe din Limitations, "The reported metric cannot
%      be decomposed" si "No score on the target labels is reported at all",
%      SCOASE: devenise false.
%      COMPOZITUL RECALCULAT NU INTRA IN TEXT, deliberat. Diferenta de 0.75
%      nu e explicabila: nu e fisierul gresit (submission.csv scoreaza
%      identic cu output/et_pca_w20.csv, adica e chiar configuratia
%      w_D = 0.20), nu e coloana gresita de gold, nu e mediere per-text
%      (59.85), nu e rotunjire (69.00), si NU e offset fix, pentru ca
%      recalculand celelalte rinduri ale ablatiei diferenta isi schimba
%      semnul (surface 10.27 vs 9.64, ET-baseline 41.76 vs 42.98, surprisal
%      64.81 vs 64.17, lemma 68.90 vs 67.89). Cauza probabila: fisierele din
%      output/ sunt re-rulari, nu exact ce s-a trimis atunci. Niciunul din
%      cele 88 de fisiere de predictie de pe disc nu ajunge la 69.74, maximul
%      e 69.18. In text intra deci compozitele returnate ca atare, plus r si
%      R^2 masurate de noi, si cele doua nu se amesteca.
%      ATENTIE, caveat onest despre r, NEscris in text: cele trei ponderi de
%      fuziune au fost alese pe scorul public, care CONTINE |r|, deci 0.777
%      nu e complet necontaminat. Sweep-ul se intinde insa pe 0.84 puncte de
%      compozit, deci cel mult 0.84/50 = 0.017 din miscare poate fi |r|.
%      Calibrarea iese confirmata, nu suspecta: R^2 e maximizat la
%      sigma = r * sigma_gold = 0.1301, iar noi am folosit 0.125, deci
%      R^2 optim ar fi fost 0.6038 fata de 0.6028 obtinut. La 0.001 de
%      optim. Iar mu = 0.27 cade pe media reala 0.2701, exact ce produce
%      optimizarea lui R^2 pe scorul public.
%   3. Kuvshynova et al. 2026 (ETRA), necitata. Combina RoLCP cu eye-tracking
%      romanesc MultiplEYE, adica exact unghiul nostru. Numele apare si in
%      Acknowledgments-ul lui RALS, deci legatura cu grupul e reala.
%   4. CALIBRAREA fata de distributia reala: REZOLVAT, nu e nimic de reparat.
%      Noi folosim mu = 0.27, sigma = 0.125 in ecuatia 2. Medie 0.23 si std
%      0.23 din Tabelul 1 al RALS descriu coloana complexity_b din RoLCP.csv,
%      NU coloana complexity, care e tinta reala si are medie 0.2701, std
%      0.1674. Verificat scorind pe fiecare coloana: complexity da 68.99,
%      complexity_b doar 56.64. Deci mu cade pe media reala si nu se comenta
%      nimic in Limitations.
%      Tot din Tabelul 2 al lor, verificat in PDF: raporteaza NUMAI Pearson,
%      niciun R^2. Diagonala RoLCP 0.71 +/- 0.05, HT 0.56 +/- .11, WT
%      0.68 +/- .04, transfer WT->RoLCP 0.60 si RoLCP->WT 0.61.
%   5. Nimic despre apendice: lucrarea nu are. Daca se adauga vreunul,
%      formatting.md cere sa vina DUPA References si cu titlu informativ,
%      iar CFP-ul MRL e autoritatea pe faptul daca sunt acceptate.
% =====================================================================

\begin{filecontents*}[overwrite]{custom.bib}
% Use this file for citations not found in the ACL Anthology (contained in "anthology.bib").

@book{Aho:72,
    author  = {Alfred V. Aho and Jeffrey D. Ullman},
    title   = {The Theory of Parsing, Translation and Compiling},
    year    = "1972",
    volume  = "1",
    publisher = {Prentice-Hall},
    address = {Englewood Cliffs, NJ}
}

@book{APA:83,
    author  = {{American Psychological Association}},
    title   = {Publications Manual},
    year    = "1983",
    publisher = {American Psychological Association},
    address = {Washington, DC}
}

@article{Chandra:81,
	author = {Ashok K. Chandra and Dexter C. Kozen and Larry J. Stockmeyer},
	year = "1981",
	title = {Alternation},
	journal = {Journal of the Association for Computing Machinery},
	volume = "28",
	number = "1",
	pages = "114--133",
	doi = "10.1145/322234.322243",
}

@inproceedings{andrew2007scalable,
  title={Scalable training of {L1}-regularized log-linear models},
  author={Andrew, Galen and Gao, Jianfeng},
  booktitle={Proceedings of the 24th International Conference on Machine Learning},
  pages={33--40},
  year={2007},
}

@book{Gusfield:97,
    author  = {Dan Gusfield},
    title   = {Algorithms on Strings, Trees and Sequences},
    year    = "1997",
    publisher = {Cambridge University Press},
    address = {Cambridge, UK}
}

@article{rasooli-tetrault-2015,
    author    = {Mohammad Sadegh Rasooli and Joel R. Tetreault},
    title     = {Yara Parser: {A} Fast and Accurate Dependency Parser},
    journal   = {Computing Research Repository},
    volume    = {arXiv:1503.06733},
    year      = {2015},
    url       = {http://arxiv.org/abs/1503.06733},
    note    = {version 2}
}

@article{Ando2005,
	Acmid = {1194905},
	Author = {Ando, Rie Kubota and Zhang, Tong},
	Issn = {1532-4435},
	Issue_Date = {12/1/2005},
	Journal = {Journal of Machine Learning Research},
	Month = dec,
	Numpages = {37},
	Pages = {1817--1853},
	Publisher = {JMLR.org},
	Title = {A Framework for Learning Predictive Structures from Multiple Tasks and Unlabeled Data},
	Volume = {6},
	Year = {2005}
}

% ---- De aici incolo, cele 20 de intrari ale lucrarii. ----

% RALS, lucrarea care publica setul nostru de test. Adaugata pe 30 iulie 2026,
% dupa ce coordonatorul a confirmat ca trebuie citata. Luata de pe
% aclanthology.org/2025.emnlp-main.1603.bib, cu O SINGURA abatere:
%
%   Anthology are al doilea autor despicat prost, "Petru-Theodor, Cristea",
%   ceea ce randeaza "Cristea Petru-Theodor". Pagina de titlu a lucrarii scrie
%   "Petru Theodor Cristea", iar intrarea cristea-nisioi-2024-machine de mai jos
%   confirma numele de familie ca "Cristea, Petru". Scris deci
%   "Cristea, Petru Theodor". A doua corectie, mai mica: Anthology scrie
%   "Creanga" fara diacritica, dar lucrarea lor scrie "Creanga" cu a-breve, deci
%   e scris {\u a}, escape-ul folosit deja la {\v{S}}tajner si {\c s}. Daca se
%   reia intrarea de pe Anthology, de refacut ambele corectii.
@inproceedings{anghel-etal-2025-rals,
    title = "{RALS}: Resources and Baselines for {R}omanian Automatic Lexical Simplification",
    author = "Anghel, Fabian  and
      Cristea, Petru Theodor  and
      Creang{\u a}, Claudiu  and
      Nisioi, Sergiu",
    editor = "Christodoulopoulos, Christos  and
      Chakraborty, Tanmoy  and
      Rose, Carolyn  and
      Peng, Violet",
    booktitle = "Proceedings of the 2025 Conference on Empirical Methods in Natural Language Processing",
    month = nov,
    year = "2025",
    address = "Suzhou, China",
    publisher = "Association for Computational Linguistics",
    url = "https://aclanthology.org/2025.emnlp-main.1603/",
    doi = "10.18653/v1/2025.emnlp-main.1603",
    pages = "31481--31492",
    ISBN = "979-8-89176-332-6"
}

% Munca anterioara pe transfer cross-lingual de complexitate pentru romana,
% prin traducere automata in engleza. Necitata pina pe 30 iulie 2026. Luata de
% pe aclanthology.org/2024.bea-1.55.bib, cu Anais convertit la {\"i}, regula
% fara-Unicode din capul fisierului.
@inproceedings{cristea-nisioi-2024-machine,
    title = "Archaeology at {MLSP} 2024: Machine Translation for Lexical Complexity Prediction and Lexical Simplification",
    author = "Cristea, Petru  and
      Nisioi, Sergiu",
    editor = {Kochmar, Ekaterina  and
      Bexte, Marie  and
      Burstein, Jill  and
      Horbach, Andrea  and
      Laarmann-Quante, Ronja  and
      Tack, Ana{\"i}s  and
      Yaneva, Victoria  and
      Yuan, Zheng},
    booktitle = "Proceedings of the 19th Workshop on Innovative Use of NLP for Building Educational Applications (BEA 2024)",
    month = jun,
    year = "2024",
    address = "Mexico City, Mexico",
    publisher = "Association for Computational Linguistics",
    url = "https://aclanthology.org/2024.bea-1.55/",
    pages = "610--617"
}

@inproceedings{shardlow-etal-2021-semeval,
    title = "{S}em{E}val-2021 Task 1: Lexical Complexity Prediction",
    author = "Shardlow, Matthew  and
      Evans, Richard  and
      Paetzold, Gustavo Henrique  and
      Zampieri, Marcos",
    editor = "Palmer, Alexis  and
      Schneider, Nathan  and
      Schluter, Natalie  and
      Emerson, Guy  and
      Herbelot, Aurelie  and
      Zhu, Xiaodan",
    booktitle = "Proceedings of the 15th International Workshop on Semantic Evaluation (SemEval-2021)",
    month = aug,
    year = "2021",
    address = "Online",
    publisher = "Association for Computational Linguistics",
    url = "https://aclanthology.org/2021.semeval-1.1/",
    doi = "10.18653/v1/2021.semeval-1.1",
    pages = "1--16"
}

@inproceedings{yimam-etal-2018-report,
    title = "A Report on the Complex Word Identification Shared Task 2018",
    author = {Yimam, Seid Muhie  and
      Biemann, Chris  and
      Malmasi, Shervin  and
      Paetzold, Gustavo  and
      Specia, Lucia  and
      {\v{S}}tajner, Sanja  and
      Tack, Ana{\"i}s  and
      Zampieri, Marcos},
    editor = "Tetreault, Joel  and
      Burstein, Jill  and
      Kochmar, Ekaterina  and
      Leacock, Claudia  and
      Yannakoudakis, Helen",
    booktitle = "Proceedings of the Thirteenth Workshop on Innovative Use of {NLP} for Building Educational Applications",
    month = jun,
    year = "2018",
    address = "New Orleans, Louisiana",
    publisher = "Association for Computational Linguistics",
    url = "https://aclanthology.org/W18-0507/",
    doi = "10.18653/v1/W18-0507",
    pages = "66--78"
}

@inproceedings{shardlow-etal-2024-multilingual,
    title = "Multilingual Resources for Lexical Complexity Prediction: A Review",
    author = "Shardlow, Matthew  and
      North, Kai  and
      Zampieri, Marcos",
    editor = "Nunzio, Giorgio Maria Di  and
      Vezzani, Federica  and
      Ermakova, Liana  and
      Azarbonyad, Hosein  and
      Kamps, Jaap",
    booktitle = "Proceedings of the Workshop on DeTermIt! Evaluating Text Difficulty in a Multilingual Context @ LREC-COLING 2024",
    month = may,
    year = "2024",
    address = "Torino, Italia",
    publisher = "ELRA and ICCL",
    url = "https://aclanthology.org/2024.determit-1.5/",
    pages = "51--59"
}

@inproceedings{shardlow-etal-2024-bea,
    title = "The {BEA} 2024 Shared Task on the Multilingual Lexical Simplification Pipeline",
    author = {Shardlow, Matthew  and
      Alva-Manchego, Fernando  and
      Batista-Navarro, Riza  and
      Bott, Stefan  and
      Calderon Ramirez, Saul  and
      Cardon, R{\'e}mi  and
      Fran{\c{c}}ois, Thomas  and
      Hayakawa, Akio  and
      Horbach, Andrea  and
      H{\"u}lsing, Anna  and
      Ide, Yusuke  and
      Imperial, Joseph Marvin  and
      Nohejl, Adam  and
      North, Kai  and
      Occhipinti, Laura  and
      Rojas, Nelson Per{\'e}z  and
      Raihan, Nishat  and
      Ranasinghe, Tharindu  and
      Salazar, Martin Solis  and
      {\v{S}}tajner, Sanja  and
      Zampieri, Marcos  and
      Saggion, Horacio},
    editor = {Kochmar, Ekaterina  and
      Bexte, Marie  and
      Burstein, Jill  and
      Horbach, Andrea  and
      Laarmann-Quante, Ronja  and
      Tack, Ana{\"i}s  and
      Yaneva, Victoria  and
      Yuan, Zheng},
    booktitle = "Proceedings of the 19th Workshop on Innovative Use of NLP for Building Educational Applications (BEA 2024)",
    month = jun,
    year = "2024",
    address = "Mexico City, Mexico",
    publisher = "Association for Computational Linguistics",
    url = "https://aclanthology.org/2024.bea-1.51/",
    pages = "571--589"
}

@inproceedings{qiu-etal-2024-complex,
    title = "{C}omp{L}ex-{ZH}: A New Dataset for Lexical Complexity Prediction in {M}andarin and {C}antonese",
    author = "Qiu, Le  and
      Guo, Shanyue  and
      Wong, Tak-Sum  and
      Chersoni, Emmanuele  and
      Lee, John  and
      Huang, Chu-Ren",
    editor = "Shardlow, Matthew  and
      Saggion, Horacio  and
      Alva-Manchego, Fernando  and
      Zampieri, Marcos  and
      North, Kai  and
      {\v{S}}tajner, Sanja  and
      Stodden, Regina",
    booktitle = "Proceedings of the Third Workshop on Text Simplification, Accessibility and Readability (TSAR 2024)",
    month = nov,
    year = "2024",
    address = "Miami, Florida, USA",
    publisher = "Association for Computational Linguistics",
    url = "https://aclanthology.org/2024.tsar-1.3/",
    doi = "10.18653/v1/2024.tsar-1.3",
    pages = "20--26"
}

@inproceedings{saggion-etal-2024-lexical,
    title = "Lexical Complexity Prediction and Lexical Simplification for {C}atalan and {S}panish: Resource Creation, Quality Assessment, and Ethical Considerations",
    author = "Saggion, Horacio  and
      Bott, Stefan  and
      Szasz, Sandra  and
      P{\'e}rez, Nelson  and
      Calder{\'o}n, Sa{\'u}l  and
      Sol{\'i}s, Mart{\'i}n",
    editor = "Shardlow, Matthew  and
      Saggion, Horacio  and
      Alva-Manchego, Fernando  and
      Zampieri, Marcos  and
      North, Kai  and
      {\v{S}}tajner, Sanja  and
      Stodden, Regina",
    booktitle = "Proceedings of the Third Workshop on Text Simplification, Accessibility and Readability (TSAR 2024)",
    month = nov,
    year = "2024",
    address = "Miami, Florida, USA",
    publisher = "Association for Computational Linguistics",
    url = "https://aclanthology.org/2024.tsar-1.9/",
    doi = "10.18653/v1/2024.tsar-1.9",
    pages = "82--94"
}

@inproceedings{conneau2020xlmr,
  title        = {Unsupervised Cross-lingual Representation Learning at Scale},
  author       = {Conneau, Alexis and Khandelwal, Kartikay and Goyal, Naman and Chaudhary, Vishrav and Wenzek, Guillaume and Guzm{\'a}n, Francisco and Grave, Edouard and Ott, Myle and Zettlemoyer, Luke and Stoyanov, Veselin},
  booktitle    = {Proceedings of the 58th Annual Meeting of the Association for Computational Linguistics (ACL)},
  pages        = {8440--8451},
  year         = {2020},
  url          = {https://aclanthology.org/2020.acl-main.747}
}

@inproceedings{pires-etal-2019-multilingual,
    title = "How Multilingual is Multilingual {BERT}?",
    author = "Pires, Telmo  and
      Schlinger, Eva  and
      Garrette, Dan",
    editor = "Korhonen, Anna  and
      Traum, David  and
      M{\`a}rquez, Llu{\'i}s",
    booktitle = "Proceedings of the 57th Annual Meeting of the Association for Computational Linguistics",
    month = jul,
    year = "2019",
    address = "Florence, Italy",
    publisher = "Association for Computational Linguistics",
    url = "https://aclanthology.org/P19-1493/",
    doi = "10.18653/v1/P19-1493",
    pages = "4996--5001"
}

@inproceedings{wu-dredze-2020-languages,
    title = "Are All Languages Created Equal in Multilingual {BERT}?",
    author = "Wu, Shijie  and
      Dredze, Mark",
    editor = "Gella, Spandana  and
      Welbl, Johannes  and
      Rei, Marek  and
      Petroni, Fabio  and
      Lewis, Patrick  and
      Strubell, Emma  and
      Seo, Minjoon  and
      Hajishirzi, Hannaneh",
    booktitle = "Proceedings of the 5th Workshop on Representation Learning for NLP",
    month = jul,
    year = "2020",
    address = "Online",
    publisher = "Association for Computational Linguistics",
    url = "https://aclanthology.org/2020.repl4nlp-1.16/",
    doi = "10.18653/v1/2020.repl4nlp-1.16",
    pages = "120--130"
}

@inproceedings{peters-etal-2019-tune,
    title = "To Tune or Not to Tune? Adapting Pretrained Representations to Diverse Tasks",
    author = "Peters, Matthew E.  and
      Ruder, Sebastian  and
      Smith, Noah A.",
    editor = "Augenstein, Isabelle  and
      Gella, Spandana  and
      Ruder, Sebastian  and
      Kann, Katharina  and
      Can, Burcu  and
      Welbl, Johannes  and
      Conneau, Alexis  and
      Ren, Xiang  and
      Rei, Marek",
    booktitle = "Proceedings of the 4th Workshop on Representation Learning for NLP (RepL4NLP-2019)",
    month = aug,
    year = "2019",
    address = "Florence, Italy",
    publisher = "Association for Computational Linguistics",
    url = "https://aclanthology.org/W19-4302/",
    doi = "10.18653/v1/W19-4302",
    pages = "7--14"
}

@inproceedings{hale2001probabilistic,
  title        = {A Probabilistic {Earley} Parser as a Psycholinguistic Model},
  author       = {Hale, John},
  booktitle    = {Proceedings of the Second Meeting of the North American Chapter of the Association for Computational Linguistics (NAACL)},
  year         = {2001},
  url          = {https://aclanthology.org/N01-1021}
}

@article{rayner1998eye,
  title        = {Eye movements in reading and information processing: 20 years of research},
  author       = {Rayner, Keith},
  journal      = {Psychological Bulletin},
  volume       = {124},
  number       = {3},
  pages        = {372--422},
  year         = {1998},
  doi          = {10.1037/0033-2909.124.3.372}
}

@inproceedings{hollenstein-etal-2021-multilingual,
    title = "Multilingual Language Models Predict Human Reading Behavior",
    author = {Hollenstein, Nora  and
      Pirovano, Federico  and
      Zhang, Ce  and
      J{\"a}ger, Lena  and
      Beinborn, Lisa},
    editor = "Toutanova, Kristina  and
      Rumshisky, Anna  and
      Zettlemoyer, Luke  and
      Hakkani-Tur, Dilek  and
      Beltagy, Iz  and
      Bethard, Steven  and
      Cotterell, Ryan  and
      Chakraborty, Tanmoy  and
      Zhou, Yichao",
    booktitle = "Proceedings of the 2021 Conference of the North American Chapter of the Association for Computational Linguistics: Human Language Technologies",
    month = jun,
    year = "2021",
    address = "Online",
    publisher = "Association for Computational Linguistics",
    url = "https://aclanthology.org/2021.naacl-main.10/",
    doi = "10.18653/v1/2021.naacl-main.10",
    pages = "106--123"
}

@inproceedings{alves-2025-benchmarking,
    title = "Benchmarking Language Model Surprisal for Eye-Tracking Predictions in {B}razilian {P}ortuguese",
    author = "Alves, Diego",
    editor = "Acarturk, Cengiz  and
      Nasir, Jamal  and
      Can, Burcu  and
      Coltekin, Cagr{\i}",
    booktitle = "Proceedings of the First International Workshop on Gaze Data and Natural Language Processing",
    month = sep,
    year = "2025",
    address = "Varna, Bulgaria",
    publisher = "INCOMA Ltd., Shoumen, BULGARIA",
    url = "https://aclanthology.org/2025.gaze4nlp-1.2/",
    pages = "7--17"
}

@inproceedings{gorman-bedrick-2019-need,
    title = "We Need to Talk about Standard Splits",
    author = "Gorman, Kyle  and
      Bedrick, Steven",
    editor = "Korhonen, Anna  and
      Traum, David  and
      M{\`a}rquez, Llu{\'i}s",
    booktitle = "Proceedings of the 57th Annual Meeting of the Association for Computational Linguistics",
    month = jul,
    year = "2019",
    address = "Florence, Italy",
    publisher = "Association for Computational Linguistics",
    url = "https://aclanthology.org/P19-1267/",
    doi = "10.18653/v1/P19-1267",
    pages = "2786--2791"
}

% Urmatoarele trei nu sunt citate in fisierul asta. Sunt citate in
% paste_body.tex (SHAP, chain-of-thought) si erau citate in subsectiunea
% 2.5 Romanian Resources, scoasa la revizia 7. Lasate aici ca sa nu mai
% fie atins nimic la asamblarea lucrarii complete.

@inproceedings{lundberg2017shap,
  title        = {A Unified Approach to Interpreting Model Predictions},
  author       = {Lundberg, Scott M. and Lee, Su-In},
  booktitle    = {Advances in Neural Information Processing Systems (NeurIPS)},
  pages        = {4765--4774},
  year         = {2017},
  url          = {https://arxiv.org/abs/1705.07874}
}

@inproceedings{wei2022cot,
  title        = {Chain-of-Thought Prompting Elicits Reasoning in Large Language Models},
  author       = {Wei, Jason and Wang, Xuezhi and Schuurmans, Dale and Bosma, Maarten and Ichter, Brian and Xia, Fei and Chi, Ed H. and Le, Quoc V. and Zhou, Denny},
  booktitle    = {Advances in Neural Information Processing Systems (NeurIPS)},
  year         = {2022},
  url          = {https://arxiv.org/abs/2201.11903}
}

@inproceedings{masala-etal-2024-vorbesti,
    title = "``Vorbe{\c s}ti Rom{\^a}ne{\c s}te?'' A Recipe to Train Powerful {R}omanian {LLM}s with {E}nglish Instructions",
    author = "Masala, Mihai  and
      Ilie-Ablachim, Denis  and
      Dima, Alexandru  and
      Corlatescu, Dragos Georgian  and
      Zavelca, Miruna-Andreea  and
      Olaru, Ovio  and
      Terian, Simina-Maria  and
      Terian, Andrei  and
      Leordeanu, Marius  and
      Velicu, Horia  and
      Popescu, Marius  and
      Dascalu, Mihai  and
      Rebedea, Traian",
    editor = "Al-Onaizan, Yaser  and
      Bansal, Mohit  and
      Chen, Yun-Nung",
    booktitle = "Findings of the Association for Computational Linguistics: EMNLP 2024",
    month = nov,
    year = "2024",
    address = "Miami, Florida, USA",
    publisher = "Association for Computational Linguistics",
    url = "https://aclanthology.org/2024.findings-emnlp.681/",
    doi = "10.18653/v1/2024.findings-emnlp.681",
    pages = "11632--11647"
}
\end{filecontents*}

\documentclass[11pt]{article}

% Change "review" to "final" to generate the final (sometimes called camera-ready) version.
% Change to "preprint" to generate a non-anonymous version with page numbers.
\usepackage[review]{acl}

% Standard package includes
\usepackage{times}
\usepackage{latexsym}

% For proper rendering and hyphenation of words containing Latin characters (including in bib files)
\usepackage[T1]{fontenc}
% For Vietnamese characters
% \usepackage[T5]{fontenc}
% See https://www.latex-project.org/help/documentation/encguide.pdf for other character sets

% This assumes your files are encoded as UTF8
\usepackage[utf8]{inputenc}

% This is not strictly necessary, and may be commented out,
% but it will improve the layout of the manuscript,
% and will typically save some space.
\usepackage{microtype}

% This is also not strictly necessary, and may be commented out.
% However, it will improve the aesthetics of text in
% the typewriter font.
\usepackage{inconsolata}

%Including images in your LaTeX document requires adding
%additional package(s)
\usepackage{graphicx}

% ---------------------------------------------------------------------
% Abaterea 3: singurele adaugiri de preambul cerute de corpul lucrarii.
% NU adaugati \bibliographystyle{acl_natbib}; acl.sty il seteaza singur,
% iar duplicarea da eroare la compilare.
% ---------------------------------------------------------------------
\usepackage{booktabs}          % \toprule / \midrule / \bottomrule in the five tables
\graphicspath{{figures/}}      % beeswarm.png, fig_corr_matrix.png

% Fractiile implicite ale lui LaTeX (topfraction 0.7, textfraction 0.2)
% dedicau coloane intregi doar figurilor, cu spatiu elastic intre ele.
% Relaxate ca sa incapa text in aceeasi coloana cu doua floats.
\renewcommand{\topfraction}{0.9}
\renewcommand{\bottomfraction}{0.5}
\renewcommand{\textfraction}{0.08}
\renewcommand{\floatpagefraction}{0.8}

% If the title and author information does not fit in the area allocated, uncomment the following
%
%\setlength\titlebox{<dim>}
%
% and set <dim> to something 5cm or larger.

\title{Cross-Lingual Transfer for Romanian Lexical Complexity Prediction:
       Fusing a Multilingual Encoder, a Language Model, Surprisal, and
       Eye-Tracking}

% Author information can be set in various styles:
% For several authors from the same institution:
% \author{Author 1 \and ... \and Author n \\
%         Address line \\ ... \\ Address line}
% if the names do not fit well on one line use
%         Author 1 \\ {\bf Author 2} \\ ... \\ {\bf Author n} \\
% For authors from different institutions:
% \author{Author 1 \\ Address line \\  ... \\ Address line
%         \And  ... \And
%         Author n \\ Address line \\ ... \\ Address line}
% To start a separate ``row'' of authors use \AND, as in
% \author{Author 1 \\ Address line \\  ... \\ Address line
%         \AND
%         Author 2 \\ Address line \\ ... \\ Address line \And
%         Author 3 \\ Address line \\ ... \\ Address line}

% ---------------------------------------------------------------------
% Blocul de autori e GOL, intentionat. MRL e double-blind.
%
% Anonimizarea o face template-ul, nu noi: acl.sty pune \acl@anonymizetrue
% pentru optiunea "review" si tipareste "Anonymous ACL submission" in locul
% lui \@author. Deci la review continutul lui \author nu se vede oricum.
% Gol e totusi mai bine decit placeholderii template-ului: daca cineva
% comuta pe "final", varianta cu "First Author / Second Author" ar pleca
% asa la tipar, iar una goala forteaza completarea.
%
% DE COMPLETAT la camera-ready, dupa modelul template-ului:
%   \author{Nume Autor \\
%     Afiliere \\
%     \texttt{email@domain} \\\And
%     Al Doilea Autor \\
%     Afiliere \\
%     \texttt{email@domain} \\}
% ---------------------------------------------------------------------
\author{}

%\author{
%  \textbf{First Author\textsuperscript{1}},
%  \textbf{Second Author\textsuperscript{1,2}},
%  \textbf{Third T. Author\textsuperscript{1}},
%  \textbf{Fourth Author\textsuperscript{1}},
%\\
%  \textbf{Fifth Author\textsuperscript{1,2}},
%  \textbf{Sixth Author\textsuperscript{1}},
%  \textbf{Seventh Author\textsuperscript{1}},
%  \textbf{Eighth Author \textsuperscript{1,2,3,4}},
%\\
%  \textbf{Ninth Author\textsuperscript{1}},
%  \textbf{Tenth Author\textsuperscript{1}},
%  \textbf{Eleventh E. Author\textsuperscript{1,2,3,4,5}},
%  \textbf{Twelfth Author\textsuperscript{1}},
%\\
%  \textbf{Thirteenth Author\textsuperscript{3}},
%  \textbf{Fourteenth F. Author\textsuperscript{2,4}},
%  \textbf{Fifteenth Author\textsuperscript{1}},
%  \textbf{Sixteenth Author\textsuperscript{1}},
%\\
%  \textbf{Seventeenth S. Author\textsuperscript{4,5}},
%  \textbf{Eighteenth Author\textsuperscript{3,4}},
%  \textbf{Nineteenth N. Author\textsuperscript{2,5}},
%  \textbf{Twentieth Author\textsuperscript{1}}
%\\
%\\
%  \textsuperscript{1}Affiliation 1,
%  \textsuperscript{2}Affiliation 2,
%  \textsuperscript{3}Affiliation 3,
%  \textsuperscript{4}Affiliation 4,
%  \textsuperscript{5}Affiliation 5
%\\
%  \small{
%    \textbf{Correspondence:} \href{mailto:email@domain}{email@domain}
%  }
%}

\begin{document}
\maketitle
\begin{abstract}
Lexical complexity prediction is normally solved by fine-tuning an
encoder on annotated data in the target language, which excludes every
language that has none. We predict the lexical complexity of Romanian
words in context without using a single Romanian training label, by
fusing four signals that carry complexity knowledge from outside the
language: cross-lingual transfer through a frozen multilingual encoder
trained on English and Romance annotations, a complexity rating and
token surprisal from a Romanian language model, and a reading-effort
score derived from eye-tracking. Together they score $69.74$ on a
combined correlation and $R^2$ metric, against $10.27$ for surface
features alone. The largest single gain
($+4.09$) comes from a lemma-frequency correction that repairs the
frequency estimates inflection corrupts, and transfers to any
morphologically rich language served by the same frequency lists. We
also report two negative results: with no target labels, fine-tuning
the encoder is consistently worse than freezing it, and in-source
validation does not predict transfer quality.
\end{abstract}

\section{Introduction}
\label{section:intro}

Making a text usable by a given audience depends, before anything else,
on its vocabulary. A word that is routine for a specialist can stop a
language learner or a reader from outside the field, taking the
surrounding sentence with it. Deciding which
words will cause this, and for whom, is what lexical complexity
prediction (LCP) does: it assigns a word in context a score on a
continuous scale \citep{shardlow-etal-2021-semeval}. Text simplification
uses those scores to choose what to replace, readability assessment to
aggregate difficulty over a document, and language-teaching material to
keep vocabulary within a learner's range; complexity prediction is for
this reason typically the first stage of a simplification pipeline
\citep{qiu-etal-2024-complex}.

Nearly every LCP system is built the same way: a pretrained encoder is
fine-tuned on complexity annotations in the language of interest. The
recipe works wherever those annotations exist, and it quietly assumes the
one resource most languages do not have. Surveying the datasets published
between 2018 and 2024, \citet{shardlow-etal-2024-multilingual} find nine
of them, covering eight languages. A language outside that set has
nothing to fine-tune on, and buying its way in means paying for a new
annotation effort of its own.

Our target language is Romanian, and we use no Romanian complexity
labels at any point, whether for training, validation or model
selection. Romanian complexity annotations do exist
\citep{anghel-etal-2025-rals}; we set them aside deliberately, to ask
what is achievable for a language that has none. Working without them
means no design decision can be checked
on Romanian itself, so every choice has to be justified on other
languages. Romanian adds a problem of its own: the most reliable
predictor of word difficulty misfires on it. Word frequency governs how
long readers fixate a word \citep{rayner1998eye} and ranks first in our
own attribution analysis (Section~\ref{sec:results}), but Romanian
inflects heavily and the frequency lists available for it cover
inflected forms poorly: of the 1587 Romanian target tokens in our test
set, $11.7\%$ are absent from the list and would be scored as maximally
rare, and for a further $36.0\%$ the listed form is rarer than its own
lemma. A method for Romanian therefore needs complexity knowledge
imported from other languages, and a frequency estimate that survives
inflection.

In this paper, we predict the lexical complexity of Romanian words in
context without a single Romanian training label, and we make four
contributions. First, we combine four label-free signals into a single
predictor: cross-lingual transfer through a frozen multilingual encoder,
a language-model complexity rating, surprisal, and a reading-effort
score derived from eye-tracking. It scores $69.74$ on our test set,
under a metric that combines Pearson correlation with $R^2$.
Second, we introduce a lemma-frequency correction for
morphologically rich languages, substituting the lemma frequency when the
inflected form is missing or rarer; it is worth $+4.09$, the largest
single gain we measure, and applies to any inflected language served
by the same frequency lists. Third, we show that the four signals are
complementary rather than redundant: the cross-lingual signal acts mainly
as a frequency detector, the rating and surprisal are near-independent,
and eye-tracking captures a difficulty the others miss, the effort of
fitting a word into its context, which shows up only in \emph{late} gaze
measures.
Fourth, we report two negative results: fine-tuning the encoder, full or
LoRA, is consistently \emph{worse} than leaving it frozen, and in-source
validation is an unreliable guide, since three configurations that win on
English-to-Romance transfer lose on Romanian.

Our setting differs from the LCP literature in the resource it assumes
rather than in the model it uses: where that literature trains on
target-language annotations, we treat complexity as knowledge to be
imported. We also invert two conventions: gaze becomes evidence about
complexity rather than something to predict from text, and the encoder
stays frozen rather than fine-tuned. Section~\ref{sec:related} sets all
three differences against the work they concern.

The rest of the paper describes the four signals and their fusion
(Section~\ref{section:approach}), reports the ablation and the analysis
(Section~\ref{sec:results}), and collects what did not work
(Section~\ref{sec:negative}).

\section{Related Work}
\label{sec:related}

\paragraph{Lexical complexity prediction.}
Annotated complexity data exists for only a handful of languages.
Romanian joined them recently: \citet{anghel-etal-2025-rals} release
$3921$ complexity-annotated Romanian samples, whose largest subset is the
test set we predict, together with a supervised baseline on it: a Ridge
regressor over handcrafted features, which reaches Pearson
$r = 0.71 \pm 0.05$ on the Romanian set under ten-fold cross-validation
grouped so that no test sentence appears in training. They report
correlation only, not $R^2$. The task
itself began as classification over
annotated target words \citep{yimam-etal-2018-report}; SemEval-2021
Task~1 replaced it with a continuous scale on the CompLex corpus
\citep{shardlow-etal-2021-semeval}, which is the formulation we adopt.
Zero-shot evaluation is not itself new: the 2018 task already had a
multilingual track whose French test set came with no French training
data, though it asked only for binary classification and supplied
training data in three languages, one of them Romance. Coverage has since
grown to
Catalan, Spanish, Mandarin and Cantonese
\citep{saggion-etal-2024-lexical, qiu-etal-2024-complex} and to the ten
languages of the MLSP shared task \citep{shardlow-etal-2024-bea}, but the
review of \citet{shardlow-etal-2024-multilingual} shows annotated LCP
data still concentrated in a few high-resource languages. Romanian
appears in none of these multilingual resources, which is why the source
side of our transfer is built from their English and Romance portions.

\paragraph{Our result on the same set.}
We predict that same set without fitting on any of its labels, and reach
Pearson $r = 0.777$, $95\%$ CI $[0.757, 0.796]$, with $R^2 = 0.603$. We
read the labels only after the predictor was frozen, in order to score
it.

\paragraph{Cross-lingual transfer without target labels.}
Multilingual encoders such as XLM-R \citep{conneau2020xlmr} place
translation-equivalent words close together, which makes zero-shot
transfer of a regressor trained on one language plausible. Transfer
quality is far from uniform across languages, however, and degrades for
languages with less pretraining data and richer morphology
\citep{pires-etal-2019-multilingual, wu-dredze-2020-languages}, which is
precisely the regime Romanian occupies. Transfer is moreover normally
realised by fine-tuning the encoder on the source language; that feature
extraction can outperform fine-tuning when the target task is distant
from the source is known \citep{peters-etal-2019-tune}, but we are aware
of no prior report of it for cross-lingual complexity transfer. This
motivates our frozen-encoder design, and the lemma-frequency correction
completes it: the encoder supplies cross-lingual semantics, the
correction repairs the lexical feature that morphology corrupts. The one
prior attempt at cross-lingual complexity prediction for Romanian takes
the translation route: \citet{cristea-nisioi-2024-machine}
machine-translate into English and predict there with English resources.
We transfer inside the shared representation space instead, which keeps
the Romanian form and its context.

\paragraph{Surprisal and eye movements as difficulty signals.}
Surprisal theory predicts that processing effort scales with the
negative log probability of a word in context \citep{hale2001probabilistic},
and gaze measures such as longer fixations, re-reading and regressions
are the standard behavioural correlate of that effort
\citep{rayner1998eye}. Language-model surprisal predicts reading times
across languages \citep{hollenstein-etal-2021-multilingual}, including
in Romance languages \citep{alves-2025-benchmarking}. Both literatures
run from text to behaviour; we take the converse view, reading surprisal
and gaze as label-free evidence about complexity.

\paragraph{Evaluation without a development set.}
Selecting a model against the same data used to report its score inflates
that score \citep{gorman-bedrick-2019-need}. With no development set to
tune on, we state our selection procedure explicitly and quantify its
effect (Section~\ref{sec:tuning}).

\section{Approach}
\label{section:approach}

\subsection{Task, data and metric}

Given a target word in a Romanian sentence, the task is to predict a
complexity score in $[0,1]$, where $0$ is a fully familiar word and $1$
a maximally difficult one. Our test data is the Romanian lexical
complexity set released with RALS \citep{anghel-etal-2025-rals} under a
CC BY-NC-SA 4.0 licence: $1587$ target words in $274$ sentences taken
from Wikipedia, popular science, literature, institutional documents and
argumentative essays, each word rated by ten annotators drawn from a pool
of eighty. Aggregated eye-tracking measurements over the same texts are
released with the complexity data, so gaze is available for the target
words without any annotation effort of our own.

Predictions are scored by
\begin{equation}
  \label{eq:metric}
  \mathrm{score} = 50\,\bigl(|r| + \max(0, R^2)\bigr),
\end{equation}
where $r$ is the Pearson correlation with the gold scores and $R^2$ the
coefficient of determination. The metric is fixed by the evaluation
setup we work in rather than chosen by us; it is not a standard LCP
measure, and it conflates two different things, since $|r|$ rewards
ranking the words correctly while $R^2$ rewards putting them on the
right absolute scale. We report it because the evaluation we submitted to
returns only this composite. The gold labels for the full test set are
public \citep{anghel-etal-2025-rals}. We used none of them to build the
predictor, and read them only once it was frozen, in order to score it.
That is what lets us report the two terms of Equation~\ref{eq:metric}
separately, $r = 0.777$ and $R^2 = 0.603$. We take $r$ as the primary
measure, for the reason given in Section~\ref{sec:tuning}.

\subsection{Setup and in-source validation}

We use XLM-RoBERTa-base \citep{conneau2020xlmr} and RoLlama3-8B
\citep{masala-etal-2024-vorbesti} in 4-bit quantisation, with
\texttt{transformers}, \texttt{scikit-learn} and \texttt{wordfreq}.

With no Romanian labels to validate on, design decisions were made
\emph{in-source}: a regressor is trained on one set of source languages
and evaluated on another, with cross-validation inside each source. We
report two such quantities, and they are not interchangeable.
$r_{\mathrm{EN}\rightarrow\mathrm{RO}^{*}}$ is the correlation obtained
when training on English CompLex \citep{shardlow-etal-2021-semeval} and
testing on the Romance portion of MLSP \citep{shardlow-etal-2024-bea};
$\bar{r}_{\mathrm{transfer}}$ averages that with the reverse direction.
Both are transfer measures, and both sit well below what the same
regressor reaches without leaving English, $r = 0.753$ on held-out
CompLex against $0.515$ transferring to Romance
(Table~\ref{tab:fail}), which is the price of crossing languages. Built
on thousands of real labels, either check tells us whether a
configuration is sound before it is applied to Romanian. Its limits are
the subject of Section~\ref{sec:negative}.

\subsection{The four signals}

Each signal estimates complexity in a different way, and none of them
requires a Romanian complexity annotation.

\paragraph{(A) Cross-lingual transfer.}
Complexity is learned from English and Romance annotations
\citep{shardlow-etal-2021-semeval, shardlow-etal-2024-bea} and carried
over to Romanian through the shared representation space of
XLM-RoBERTa, which places translation-equivalent words close together
(Section~\ref{sec:related}). A Ridge and HistGradientBoosting blend
reads each word's embedding, the mean of the last four layers centred
per language, together with a few surface cues: frequency, length,
syllable count and capitalisation.

One Romanian-specific correction accounts for the largest single gain we
measure. Because Romanian is heavily inflected, many word forms are
absent from the frequency list and would be scored as maximally rare: of
the $1587$ target tokens, $11.7\%$ have no entry at all, and for
$36.0\%$ the lemma is more frequent than the inflected form, by $0.86$
Zipf units on average. The damage is visible in the annotations
themselves: \citet{anghel-etal-2025-rals} measure a correlation of
$-0.53$ between Zipf frequency and Romanian complexity, against $-0.74$
for the English data annotated under the same guidelines. We therefore
use
$\max\bigl(\mathrm{zipf}(\mathrm{form}), \mathrm{zipf}(\mathrm{lemma})\bigr)$
in place of the form frequency.

\paragraph{A scope condition for the correction.}
\label{sec:scope}
The substitution assumes that an inflected form inherits its difficulty
from its lemma, and that assumption has a failure mode worth naming.
\emph{Putu}, the simple-perfect form of \emph{putea} (``to be able''), is
rare, archaic in ordinary prose, and hard for exactly that reason; its
lemma is among the most frequent verbs in the language. The correction
raises its frequency estimate by $2.76$ Zipf units and so predicts it to
be easy. The assumption fails whenever the inflection is itself the
source of difficulty, which is the case for archaic, literary or
otherwise marked morphology. The $+4.09$ measured over $1587$ items says
the assumption usually holds; forms like this one say it is an
assumption, not an identity.

\paragraph{(B) Language-model rating.}
RoLlama3 is asked to rate the target word in context on a $1$--$5$
scale. We read the answer from the probabilities the model assigns to
the five digit tokens rather than from the text it generates, which
turns a discrete choice into a smooth expected value.

\paragraph{(C) Surprisal.}
Surprisal, $-\log P(w \mid \mathrm{context})$, measures how unexpected
the word is; under surprisal theory an unexpected word is harder to
process \citep{hale2001probabilistic}. It comes from the same forward
pass as (B) but is a raw internal statistic, which is why it is nearly
independent of the rating ($r = 0.07$).

\paragraph{(D) Eye-tracking.}
Harder words are fixated longer and re-read more often
\citep{rayner1998eye}. We compress the eight per-word effort measures
into a single reading-effort score with PCA. The first component
explains $62\%$ of the variance and every measure loads on it with the
same sign (Table~\ref{tab:pca}), so the component is interpretable as a
single ``more effort $=$ more complex'' axis.

\begin{table}[tbp]
\centering\small
\begin{tabular}{lc}
\toprule
\textbf{Eye-tracking metric} & \textbf{PC1 loading} \\
\midrule
Total fixation time (TFT) & 0.44 \\
Total fixation count (TFC) & 0.43 \\
First-pass reading time (FPRT) & 0.42 \\
First-pass fixation count (FPFC) & 0.39 \\
Re-reading time (RRT) & 0.35 \\
First-fixation duration (FFD) & 0.29 \\
Regression-path duration (RPD\_inc) & 0.28 \\
Incoming regressions (TRC\_in) & 0.13 \\
\bottomrule
\end{tabular}
\caption{Weight of each eye-tracking metric in the reading-effort score (loadings on PC1, which explains $62\%$ of
the variance). All loadings are positive: one ``more effort = more complex'' axis.}
\label{tab:pca}
\end{table}

\subsection{Fusing the signals}

The four signals live on different scales, so each is standardised and
they are combined as a weighted average. Signal A is itself an ensemble
of the English-source and Romance-source transfers, averaged after
standardisation. The fused score is then mapped onto the scale of the
test distribution,
\begin{equation}
  \label{eq:calib}
  \hat{y} = \mathrm{clip}\bigl(\mu + \sigma \cdot z(s),\; 0,\; 1\bigr),
\end{equation}
with $\mu = 0.27$ and $\sigma = 0.125$. Per-signal weights are in
Table~\ref{tab:signals}.

\begin{table}[tbp]
\centering\small
\begin{tabular}{lcc}
\toprule
\textbf{Signal} & \textbf{SHAP importance} & \textbf{weight} \\
\midrule
A cross-lingual & 0.061 & 44\% \\
B LLM rating & 0.034 & 29\% \\
D eye-tracking & 0.017 & 20\% \\
C surprisal & 0.006 & 6\% \\
\bottomrule
\end{tabular}
\caption{Per-signal SHAP importance and the weight each signal gets in the final average. The two orders agree;
SHAP describes the model, it does not choose the weights.}
\label{tab:signals}
\end{table}

\subsection{What was tuned, and what it costs}
\label{sec:tuning}

A zero-label setting has no development set, and selecting anything
against the evaluation score is a known way to inflate it
\citep{gorman-bedrick-2019-need}. We state exactly what this affects.

Five numbers were chosen against the public evaluation score: the three
fusion weights and the two constants of Equation~\ref{eq:calib}. The
weights were searched on a coarse grid; for the eye-tracking weight the
full sweep was
$w_D \in \{0, 0.10, 0.20, 0.30\} \rightarrow 68.90, 69.66, 69.74,
69.00$. Everything else, the encoder, the layer, the feature set, the
regressor and the lemma correction, was decided in-source, without
reference to that score.

Two things bound the damage. First, the calibration constants cannot
affect $r$ at all. Equation~\ref{eq:calib} is affine up to the clip,
and the clip is inactive on our predictions, whose range is
$[0.014, 0.845]$; Pearson $r$ is invariant under affine rescaling, so
$\mu$ and $\sigma$, the two quantities that most directly encode
information about the test distribution, move only the $R^2$ term of
Equation~\ref{eq:metric}. This is why we treat $r$ as the primary
measure. Second, the response surface is flat: the eye-tracking sweep
spans $0.84$ points end to end, so no setting of $w_D$ within the grid
changes the outcome by much. The one figure that inherits the full
optimism of this procedure is the $+0.84$ credited to eye-tracking in
Table~\ref{tab:ablation}, flagged there as an upper bound.

\begin{table}[tbp]
\centering\small
\begin{tabular}{lcc}
\toprule
\textbf{Configuration} & \textbf{Score} & \textbf{$\Delta$} \\
\midrule
Predict a constant (the mean) & 0.00 & --- \\
Surface features only (zero-shot RO) & 10.27 & +10.27 \\
\;+ eye-tracking (no imported res.) & 41.76 & +31.49 \\
\;+ cross-lingual A + LLM B & 64.52 & +22.76 \\
\;+ surprisal C & 64.81 & +0.29 \\
\;+ lemma-frequency fix (in A) & 68.90 & +4.09 \\
\;+ eye-tracking PCA D & \textbf{69.74} & +0.84$^{\dagger}$ \\
\bottomrule
\end{tabular}
\caption{Evaluation score (Equation~\ref{eq:metric}), built one component at a time; $\Delta$ is the marginal
contribution of each component. Predicting one constant gives no correlation, hence $0$. The first two rows use
only resources released with the data. $^{\dagger}$This delta is the maximum of the sweep that selected $w_D$ and
should be read as an upper bound (Section~\ref{sec:tuning}).}
\label{tab:ablation}
\end{table}

\section{Results}
\label{sec:results}

\subsection{How much each signal contributes}

Table~\ref{tab:ablation} builds our predictor one signal at a time.
Restricted to the resources released with the data, surface features and
eye-tracking, it reaches $41.76$; the imported signals take it to
$69.74$, so roughly $40\%$ of the final score is attributable to
knowledge brought in from outside Romanian. Surface features alone reach
$10.27$. The single largest step is adding the cross-lingual and
language-model signals together ($+22.76$); the lemma-frequency
correction is the largest single gain after that ($+4.09$), and
surprisal and eye-tracking each add a small amount.

Against the gold labels our predictions reach $r = 0.777$, $95\%$ CI
$[0.757, 0.796]$, and $R^2 = 0.603$ over all $1587$ target words.
\citet{anghel-etal-2025-rals} report $r = 0.71 \pm 0.05$ on the same
words with a Ridge regressor fitted on them under grouped ten-fold
cross-validation. Theirs measures what a small supervised model extracts
from in-domain Romanian annotation, ours what a frozen fusion imports
from outside the language, and their figure remains the right measure of
what target-language annotation buys.

\begin{figure}[tbp]
\centering
\includegraphics[width=\columnwidth]{beeswarm.png}
\caption{SHAP over named features (top 8 shown). Frequency dominates, then the LLM rating and surprisal. The
separate blue cluster on the frequency row is the group of very rare / out-of-vocabulary words, which push the
predicted complexity up, the expected behavior. Red = high feature value.}
\label{fig:beeswarm}
\end{figure}

\subsection{What drives the prediction}

A SHAP analysis \citep{lundberg2017shap} over named features
(Figure~\ref{fig:beeswarm}) places
frequency first, the language-model rating second and surprisal third;
the surface features barely matter. Table~\ref{tab:signals} gives the
same ordering at the signal level, and it agrees with the tuned weights.
The signals are complementary rather than redundant
(Figure~\ref{fig:corr}): the rating and surprisal are nearly
independent, while the cross-lingual signal tracks frequency closely
enough to be read as a frequency detector with semantics attached.

\begin{figure}[tbp]
\centering
\includegraphics[width=\columnwidth]{fig_corr_matrix.png}
\caption{Correlation between the four signals and frequency on the Romanian test set. The rating (B) and surprisal
(C) are nearly independent; the cross-lingual signal (A) tracks frequency closely. Every cell is computed on
$n = 1587$, where any $|r| > 0.05$ already reaches $p < 0.05$, so the matrix is to be read by shared variance
rather than by significance: A shares $83\%$ of its variance with frequency, B and C share $0.5\%$.}
\label{fig:corr}
\end{figure}

\subsection{Eye-tracking contributes through its late measures}

Eye-tracking helps because it carries something the lexical signals do
not, and the division is the one predicted by the reading literature.
Partialling out frequency and length from each gaze measure over the
$1509$ test words that have gaze data, the early measures (first
fixation duration, first-pass reading time, skip rate) retain almost no
relation to our own prediction ($\bar{r} = 0.010$) or to surprisal
($\bar{r} = -0.027$): what they encode is lexical access, which signal A
already knows. The late measures (re-reading time, incoming regressions,
refixation rate) retain $\bar{r} = 0.243$ and $\bar{r} = 0.177$
respectively. Late gaze is where the contextual-integration difficulty
lives, and that is the part of the eye-tracking signal that is not
already in the other three.


\section{What Did Not Work}
\label{sec:negative}

Table~\ref{tab:fail} lists the attempts that failed, and two of them
carry a lesson beyond this work.

\paragraph{Fine-tuning is worse than freezing.}
Every form of fine-tuning we tried, full fine-tuning and LoRA over a
sweep of learning rates, epochs and ranks, transferred worse than the
frozen encoder ($0.461$ and $0.462$ against $0.515$). With no
target-language labels, fine-tuning can only adapt to the source
language, and in doing so it degrades the cross-lingual alignment that
the transfer depends on. This is the behaviour \citet{peters-etal-2019-tune}
report when the target task is distant from the source, and it is worth
stating plainly because the default move in a cross-lingual pipeline is
to fine-tune \citep{pires-etal-2019-multilingual}.

\paragraph{In-source validation does not predict transfer to Romanian.}
Three configurations beat the frozen XLM-R baseline on in-source
transfer and then lost on Romanian: a different encoder (e5), mid-layer
embeddings, and a Lasso regressor. The in-source check measures transfer
between English and Romance, related and well-resourced languages, and
that is evidently not the same problem as transfer into Romanian. In a
setting with no target-language development set this is the
uncomfortable conclusion: the validation signal available is not the
validation signal needed.

\paragraph{Prompting and extra features.}
Chain-of-thought and few-shot prompting added noise rather than signal;
at the $8$B scale chain-of-thought is not yet reliably helpful
\citep{wei2022cot}. Rarity measures and a range of part-of-speech,
named-entity and WordNet features proved redundant with frequency,
which is consistent with the SHAP picture in
Figure~\ref{fig:beeswarm}.

\begin{table}[tbp]
\centering\small
\begin{tabular}{lcc}
\toprule
\textbf{Attempt} & \textbf{in-source} & \textbf{score} \\
\midrule
\multicolumn{3}{l}{\emph{Fine-tuning} ($r_{\mathrm{EN}\rightarrow\mathrm{RO}^{*}}$)} \\
Frozen XLM-R (this work) & 0.515 & --- \\
Full fine-tune (best of sweep) & 0.461 & --- \\
LoRA (best of sweep) & 0.462 & --- \\
\midrule
\multicolumn{3}{l}{\emph{Representation and regressor} ($\bar{r}_{\mathrm{transfer}}$)} \\
Frozen XLM-R (this work) & 0.551 & 69.74 \\
Encoder swap (e5) & 0.586 & 63.00 \\
Mid-layer embeddings & 0.556 & 64.49 \\
Lasso regressor & 0.569 & 64.28 \\
\midrule
\multicolumn{3}{l}{\emph{Prompting}} \\
LLM few-shot prompt & --- & 63.31 \\
\bottomrule
\end{tabular}
\caption{Failed attempts. The two blocks use different in-source measures and are not comparable across blocks:
$r_{\mathrm{EN}\rightarrow\mathrm{RO}^{*}}$ is the one-directional English-to-Romance transfer correlation, and
$\bar{r}_{\mathrm{transfer}}$ averages both directions. Within the second block, every variant that beats the
frozen baseline in-source loses on the public evaluation score.}
\label{tab:fail}
\end{table}

\section{Conclusions}
\label{section:conclusions}

Working without Romanian training labels did not merely constrain the
design, it determined it: with no labels to fit and no development set
to select on, the only remaining move is to import complexity knowledge
from elsewhere and combine several independent imports. Four such
signals, a frozen multilingual encoder, a language-model rating,
surprisal and gaze, reach $69.74$ on the metric of
Equation~\ref{eq:metric} without a single target-language annotation.

Three findings generalise beyond Romanian. A lemma-frequency correction
is worth more than any other single component ($+4.09$) for a
morphologically rich language, subject to the scope condition of
Section~\ref{sec:scope}. Fine-tuning a multilingual encoder without
target labels is actively harmful, and freezing it is the stronger
default. And in-source validation, the only validation a zero-label
setting affords, does not reliably predict transfer to the target
language, which is the finding we would most want future work in this
setting to take seriously.

\section*{Limitations}

\paragraph{Five quantities were selected on the evaluation score.}
The setting provides no development set, so the three fusion weights
and the two calibration constants of Equation~\ref{eq:calib} were
chosen against the public evaluation score rather than held-out data. We
have bounded this in Section~\ref{sec:tuning}, but only partly: the
calibration cannot influence $r$, which is algebraic, and the response
surface is flat across the grid we searched. Neither argument measures
the bias that remains, and we have no held-out Romanian labels with
which to measure it. The $69.74$ headline figure and, in particular, the
$+0.84$ credited to eye-tracking therefore retain an optimistic bias we
cannot quantify.

\paragraph{The eye-tracking analysis is partly circular.}
The early-versus-late analysis in Section~\ref{sec:results} was not
validated against the gold labels. It uses our own predictions as the
complexity proxy over the $1509$ test words that have gaze, so it
establishes that late gaze measures carry information the other signals
lack \emph{relative to what our model predicts}, not relative to human
judgement. Gaze coverage is also not missing at random: the $78$ words
without fixations are shorter and more frequent than average (mean Zipf
$4.45$ against $4.20$; length $6.91$ against $7.57$), because frequent
short words are skipped, so the signal is systematically absent exactly
where it would matter least, and its absence is itself informative in a
way the PCA score does not capture.

\paragraph{The lemma-frequency correction has a scope condition.}
Substituting lemma frequency assumes an inflected form inherits its
difficulty from its lemma. Section~\ref{sec:scope} documents a case
where this fails, a simple-perfect verb form that is difficult because
of its inflection, and we have not measured how often such forms occur,
since doing so would require Romanian labels. The correction should be
expected to help on average and to fail systematically on archaic,
literary or otherwise marked morphology.

\paragraph{The approach assumes resources that not every language has.}
The method needs an instruction-tuned language model in the target
language, a frequency list with lemma coverage, and annotated complexity
data in a related language. Romanian has all three, and the eye-tracking
measurements came released with the complexity data, which makes gaze
the least portable of the four signals: gaze corpora exist for few
languages, and collecting them is expensive. A language lacking any of
these would be left with a proper subset of our signals, and we have not
measured how the fusion degrades under such a subset.

\paragraph{One target language, one domain.}
All results are for Romanian, on the texts of a single dataset. The
claims we expect to generalise, that a frozen encoder beats fine-tuning
without target labels and that in-source validation misleads, are
supported here by one language pair each, and neither has been tested on
a second morphologically rich language.

% ---------------------------------------------------------------------
% Blocul \section*{Acknowledgments} al template-ului a fost SCOS de aici.
% formatting.md, la Acknowledgments: "Do not include this section in the
% review version." Nu e o optiune.
%
% DE REPUS la camera-ready, cu mulitumirile noastre, si exact in locul
% asta: formatting.md cere ca sectiunea sa fie imediat inaintea
% referintelor si sa fie nenumerotata.
% ---------------------------------------------------------------------

% Bibliography entries for the entire Anthology, followed by custom entries
%\bibliography{custom,anthology-overleaf-1,anthology-overleaf-2}

% Custom bibliography entries only
\bibliography{custom}

% ---------------------------------------------------------------------
% NU exista apendice. Apendicele de exemplu al template-ului ("Example
% Appendix" / "This is an appendix.") a fost scos, iar tab:pca, mutat un
% timp aici, s-a intors in 3.3: singur intr-un apendice, dupa cele 18
% referinte, ocupa o pagina intreaga cu 14 rinduri de text, si arata rau.
% Costul intoarcerii e ~1/4 de pagina de continut, care exista.
%
% Daca se adauga apendice mai tirziu, formatting.md cere: sa vina DUPA
% References, sa aiba titlu informativ, iar textul principal sa rimina
% autonom fara ele. Si sa fie confirmate in CFP-ul MRL, care e autoritatea
% pe lungime si pe acceptarea apendicelor.
% ---------------------------------------------------------------------

\end{document}
